\chapter{A constructive recall of locally presentable categories}\label{ch:lp}

In this chapter we recall some basic theory of locally presentable categories following~\citet{adamek1994locally}. The reason we are redeveloping some of these materials is that in this document we have adopted a constructive meta-theory, and the treatment of locally presentable categories in \emph{loc.\ cit.}\ is not always evidently constructive. 

Our choice for a constructive meta-theory is \emph{not} philosophical. One technique that will be useful for our development of \gls{QC} for Horn theories with higher arity is to interpret an externally proven fact in the internal logic of a suitable topos. Such an internalisation process is valid only if the internal logic of the topos supports all the logical resources used in the original proof, and in general a topos only supports constructive logic. Thus, a constructive development of the theory of locally presentable categories, despite of interest in itself, is also quite useful for developing the subject of this document.

Though the theory of locally presentable categories is well-established, one of the innovative parts of our constructive treatment is to find a suitable notion of \emph{size}, that can replace the notion of \emph{regular cardinal} in classical set theory.

\section{A notion of size}\label{sec:size}

\begin{definition}[Size]\label{def:size}
  A \emph{notion of size} $\kappa$ is a collection of sets such that:
  \begin{itemize}
    \item $\kappa$ is \emph{replete}, i.e.\ it is closed under isomorphisms of sets;
    \item $\kappa$ contains all \emph{finite cardinals} $\mb n$ for $n\in\N$;
    \item $\kappa$ is closed under \emph{dependent sums}: for any $A \in \kappa$ and a family $B \colon A \to \kappa$, the disjoint union $\sum_{a\in A}B(a)$ belongs to $\kappa$;
    \item sets in $\kappa$ are \emph{projective};
  \end{itemize}
  We say $\kappa$ is \emph{separated}, if it is also closed under identities:
  \begin{itemize}
    \item for any $X\in\kappa$ and $x,y\in X$, the proposition $x = y$ also belongs to $\kappa$.
  \end{itemize}
  We say $\kappa$ is \emph{essentially small}, if it is classified by a small set:
  \begin{itemize}
    \item there exists a small set $K$ and a map $\El \colon K \to \kappa$, such that for any $\kappa$-small set $I$, there exists $k\in K$ such that $I \cong \El(k)$.
  \end{itemize}
  A set is \emph{$\kappa$-small} if it belongs to $\kappa$, and $\kappa\hp\Set$ denotes the full subcategory of $\kappa$-small sets.
\end{definition}

\begin{remark}[Comparison of the notion of size]\label{rem:comparethenotionofsize}
  Classically in \gls{ZFC} the above notion of an essentially small (separated) size corresponds exactly to an infinite \emph{regular cardinal}: under \gls{LEM}, being separated is automatic since every set is decidable; under choice, every set is also projective; thus the remaining condition of containing all finite sets and being closed under dependent sums exactly corresponds to an infinite cardinal. Constructively, a very similar notion of size was also independently proposed in a recent work~\citet{coquand_et_al:LIPIcs.LICS.2026.31}, though they do not require a size to be separated; this condition is to make sure $\kappa$-small sets induce $\kappa$-small discrete groupoids (cf.~\zcref{exa:discgroupoidkappasmall}).
\end{remark}

\begin{remark}[Closure properties of $\kappa$-small sets]\label{rem:closurepropofkappa}
  For a notion of size $\kappa$, being closed under dependent sums, together with the other requirements specified in~\zcref{def:size}, has many useful consequences for closure properties of the category $\kappa\hp\Set$ of $\kappa$-small sets. For instance, $\kappa\hp\Set$ will be closed under \emph{$\kappa$-small coproducts}, in particular it will be closed under \emph{cartesian products}. If $\kappa$ is \emph{separated}, then it is also closed under \emph{equalisers}, since given $f,g \colon X \to Y$, the equaliser is given by the subset 
  \[ \mr{Eq}(f,g) \cong \scomp{x\in X}{fx = gx} \cong \sum_{x\in X}fx = gx. \]
  Hence, $\kappa\hp\Set$ is closed under $\kappa$-small coproducts and finite products in $\Set$, and if $\kappa$ is separated then it is in fact closed under all finite limits. However, in general it is \emph{not} closed under coequalisers (cf.~\zcref{rem:coequalisercntandmarkov}).
\end{remark}

For our purpose, we are exclusively interested in the following two notions of size:

\begin{example}[Finiteness as a notion of size]\label{exa:finitesize}
  We say a set $X$ is \emph{finite} if it is isomorphic to a finite cardinal,
  \[ \ms{isFin}(X) \coloneq \ex n{\N} X \cong \mb n. \]
  Finite sets are \emph{decidable}, thus they are closed under identities. Finite sets are also closed under dependent sums, which can be shown by induction. By construction, finite sets are also projective. Furthermore, finite sets by definition has a classifier $\N$. Thus, finite sets form a separated and essentially small size in the sense of~\zcref{def:size}, and we denote this size as $\omega$. The full subcategory $\omega\hp\Set$ of finite sets will also be denoted as $\Fin$. 
\end{example}

\begin{example}[Countability as a notion of size]\label{exm:countablesize}
  A set $X$ is \emph{countable} if it is isomorphic to a complemented subset of the natural numbers,
  \[ \ms{isCnt}(X) \coloneq \ex{I}{\mb 2^\N} X \cong \scomp{n\in\N}{I(n)}. \]
  The collection of countable sets will be denoted as $\omega_1$. By construction, all countable sets are (isomorphic to) complemented subsets of $\N$, thus they are all decidable. This in particular means $\omega_1$ is closed under identities. By countable choice, $\N$ is projective, and thus by~\zcref{lem:projdecsubset} all countable sets are projective. Suppose $X$ is countable and $Y \colon X \to \omega_1$ is a countable family of countable sets. To show $\sum_{x:X}Y_x$ is countable, we may assume $X$ is isomorphic to $I$ for some complemented subset $I$ of $\N$. Then by assumption, 
  \[ \fa{i}{I}\ex{J}{\mb 2^\N} Y_i \cong \scomp{n\in\N}{J(n)}. \] 
  By $I$ being projective, there exists a family $J_i$ of complemented subsets of $\N$ such that $Y_i \cong J_i$. Thus, the dependent sum $\sum_{x:X}Y_x$ will be equivalent to $\sum_{i:I}J_i$, which is a complemented subset of $\N \times \N$. Via a bijection $\N \times \N \cong \N$, this shows $\sum_{x:X}Y_x$ is countable. Furthermore, again by construction $\omega_1$ has a classifier $\mb 2^\N$. This shows $\omega_1$ is also a separated essentially small size. We denote the category $\omega_1\hp\Set$ of countable sets also as $\Cnt$.
\end{example}

\begin{remark}[Coequalisers of countable sets]\label{rem:coequalisercntandmarkov}
  $\Fin$ can be shown to be a Boolean elementary topos, see e.g.~\citet[D5.2]{johnstone2002sketches}. In particular, it has coequalisers. On the other hand, constructively $\Cnt$ is not necessarily closed under coequalisers. For instance, it is not hard to see that $\Cnt$ having coequalisers implies the \emph{Markov principle}.
\end{remark}

At the end of this section, we observe that any notion of size $\kappa$ induces a \emph{dominance} in the sense of~\citet{rosolini1986continuity}:

\begin{definition}[Dominance]\label{def:dominance}
  A dominance $\Sigma$ is a subuniverse of $\pp$ that is closed under dependent sums, i.e.\ $\mb 1\in\Sigma$, and for any $\varphi\in\Sigma$ together with a family $\psi \colon \varphi \to \Sigma$, the dependent sum $\sum_{x\in\varphi}\psi(x)$ also belongs to $\Sigma$.
\end{definition}

\begin{definition}[$\kappa$-overt propositions]\label{def:kappaopen}
  A proposition is \emph{$\kappa$-overt}, if it is of the form $\pss X$ for some $\kappa$-small set $X$. The subuniverse of $\kappa$-overt propositions will be denoted as $\mb{O}_\kappa$. 
\end{definition}

\begin{proposition}\label{prop:sizedominance}
  For any notion of size $\kappa$, $\mb{O}_\kappa$ will be a dominance closed under finite joins.
\end{proposition}
\begin{proof}
  Note that $\mb 1\in\mb{O}_\kappa$. It is closed under finite joins because $\kappa$-small types are closed under finite coproduct (cf.~\zcref{rem:closurepropofkappa}). Now suppose we have $\varphi \in \mb{O}_\kappa$, and consider a family $\psi \colon \varphi \to \mb{O}_\kappa$ of $\kappa$-overt propositions indexed by $\varphi$. By definition, this means we have
  \[ \fa{i}{\varphi}\ex{Y}{\kappa}\psi_i = \pss{Y}. \]
  To show the dependent sum $\sum_{i\in\varphi}\psi_i$ belongs to $\mb{O}_\kappa$, by assumption we may assume $\varphi$ is of the form $\pss X$ for some $X \in \kappa$. In particular, there is a surjection $q \colon X \surj \varphi$. Hence,
  \[ \fa{x}{X}\ex Y{\kappa} \psi_{qx} = \pss{Y}. \]
  Since $X\in\kappa$, $X$ is projective, thus there exists an $X$-indexed family of $\kappa$-small sets $Y \colon X \to \kappa$, where $\psi_{qx} = \pss{Y_x}$ for all $x\in X$. This way, the dependent sum can be computed by
  \[ \sum_{i\in\varphi}\psi_i = \pss{\sum_{x\in X}Y_x}. \]
  Hence, $\mb{O}_\kappa$ is closed under dependent sums, and forms a dominance.
\end{proof}

\begin{remark}[Relation to the dominance of semi-decidable propositions]
  It is shown by~\citet[Prop. 5.2.1]{rosolini1986continuity} (with the proof is attributed to Dana Scott) that the following subuniverse of \emph{semi-decidable propositions} is a dominance,
  \[ \Sigma_{\mr{s.d.}} \coloneq \scomp{p\in\pp}{\ex f{\N^\N} p = \ex{n}{\N}fn = 0}. \]
  In fact, it is not hard to see that the universe $\mb O_{\omega_1}$ of $\omega_1$-overt propositions is identical to $\Sigma_{\mr{s.d.}}$. The proof in \emph{loc.\ cit.}\ also crucially relies on countable choice. In fact, the proof of~\zcref{prop:sizedominance} can be seen as a suitable generalisation of this proof that works for any notion of size. For $\omega$ we have $\mb O_\omega = \mb 2$, recovering the dominance of decidable propositions.
\end{remark}

\section{Filtered diagrams}

In this section we fix a \emph{separated size} $\kappa$. We say a category $\mc A$ is \emph{$\kappa$-small} if its set of objects is $\kappa$-small, and for any $a,b \in \mc A$, the hom-set $\mc A(a,b)$ is $\kappa$-small. 

\begin{example}[Discrete groupoids from $\kappa$-smalls sets are $\kappa$-small]\label{exa:discgroupoidkappasmall}
  The requirement for $\kappa$ to be separated is to make $\kappa$-small sets also $\kappa$-small discrete groupoids. For any $X\in\kappa$, the hom-sets in the associated discrete groupoid is given by $X(x,y) \cong (x = y)$ for any $x,y\in X$. By assumption the identity proposition $x = y$ is $\kappa$-small since $\kappa$ is separated, hence $X$ as a discrete groupoid will be a $\kappa$-small category.
\end{example}

\begin{definition}[$\kappa$-(co)completeness]
  A category $\mc A$ is $\kappa$-(co)complete if it has $\kappa$-small (co)limits, i.e.\ if for any $\kappa$-small category $\mc D$ and diagram $F \colon \mc D \to \mc A$, the (co)limit of $F$ exists.
\end{definition}

By~\zcref{exa:discgroupoidkappasmall}, diagrams indexed by $\kappa$-small sets will in particular be instances of $\kappa$-small (co)limits, thus $\kappa$-(co)complete categories have $\kappa$-small (co)products. Since any $\kappa$ also contains all finite sets, $\kappa$-(co)complete categories also have (co)equalisers. In fact, the existence of these two classes of (co)limits classifies $\kappa$-(co)completeness:

\begin{proposition}\label{prop:smallcompiffeqcoandeq}
  $\mc A$ is $\kappa$-(co)complete if it has (co)equalisers and $\kappa$-small (co)products.
\end{proposition}
\begin{proof}
  The only if direction is trivial. For the if direction, for a $\kappa$-small diagram $F \colon \mc D \to \mc A$, its limit is given by the following equaliser,
  \[ 
  \begin{tikzcd}
    \lt F \ar[r, dashed, tail] & \prod_{d\in\Ob(\mc D)}Fd \ar[r, shift left] \ar[r, shift right] & \prod_{\alpha\in{\mc D}(d_0,d_1)}Fd_1
  \end{tikzcd}
  \]
  It then suffices to observe that $\prod_{d\in\Ob(\mc D)}Fd$ and $\prod_{\alpha\in\mc D(d_0,d_1)}Fd_1$ are $\kappa$-small products: the former by assumption that $\Ob(\mc D)$ is $\kappa$-small, and the latter by the fact that the dependent sum $\sum_{d_0,d_1\in\Ob(\mc D)}\mc D(d_0,d_1)$ is $\kappa$-small, since by~\zcref{rem:closurepropofkappa} $\kappa$-small sets are closed under finite products and dependent sums.
\end{proof}

\begin{definition}[$\kappa$-filtered category]
  We say a category $\mc I$ is \emph{$\kappa$-filtered}, if it is small and every $\kappa$-small diagram $\mc D \to \mc I$ admits a cocone.
\end{definition}

\begin{remark}[$\kappa$-cocomplete categories are $\kappa$-filtered]\label{rem:kappacompeltefiltered}
  If $\mc I$ has all $\kappa$-small colimits, then in particular every $\kappa$-small diagram has a universal cocone, hence it is $\kappa$-filtered.
\end{remark}

As usual, $\kappa$-filtered categories can be characterised in a more concrete fashion:

\begin{lemma}\label{lem:concretefilter}
  A category $\mc I$ is $\kappa$-filtered if and only if the following hold:
  \begin{itemize}
    \item For any $\kappa$-small set $I$ and any diagram $f \colon I \to \mc I$, there exists $x\in\mc I$ and a family of morphisms $\alpha \in \prod_{i\in I}\mc I(fi,x)$;
    \item For any $\kappa$-small set $I$, $x,y\in\mc I$, and any $I$-indexed family of arrows $\alpha \colon I \to \mc J(x,y)$, there exists $z\in\mc J$ and $\gamma \colon y \to z$ that coequalises them, i.e.\ for all $i,j\in I$, $\gamma \circ \alpha_i = \gamma \circ \alpha_j$.
  \end{itemize}
\end{lemma}
\begin{proof}
  Necessity is evident. For sufficiency, we essentially follow the proof strategy for the finitary case; see e.g.~\citet[Lem. 2.13.2]{borceux1994handbookI}. Let $F \colon \mc D \to \mc I$ be a $\kappa$-small diagram in $\mc I$. We construct a cocone as follows: 
  \begin{itemize}
    \item First pick $k_0 \in \mc I$ with a family $\alpha \in \prod_{d\in\mc D}{\mc I}(Fd,k_0)$;
    \item For any $f \in {\mc D}(d,e)$, pick $k_f$ with $\beta_f \in {\mc I}(k_0,k_f)$, coequalising the pair $\alpha_d \circ Ff$ and $\alpha_e$;
    \item We now again have a family $k_f$ indexed by $\sum_{d,e\in\mc D}{\mc D}(d,e)$, thus we may find some $k_1$ with a family of maps $\gamma_f \in {\mc I}(k_f,k_1)$.
    \item Thus, we get a $\kappa$-small family of maps $\gamma_-\circ\beta_- \colon k_0 \to k_1$ in $\mc I$, indexed by $\sum_{d,e\in\mc D}{\mc D}(d,e)$. By assumption, we can find $k$ with $\delta \in {\mc I}(k_1,k)$ coequalising them all.
  \end{itemize}
  We then define the cocone to be consisting of the vertex $k$, with component for each $d\in\mc D$
  \[ \epsilon_d \coloneq \delta \circ \gamma_{1_d} \circ \beta_{1_d} \circ \alpha_d \in {\mc I}(Fd,k). \]
  This is a cocone by construction. 
\end{proof}

In $\Set$, colimits indexed by $\kappa$-filtered categories are easier to compute, since one can explicitly characterise the equivalence relation generating the colimit:

\begin{lemma}\label{lem:filteredquotient}
  Let $F \colon \mc I \to \Set$ be a $\kappa$-filtered diagram. Then the following relation $\sim$ defined on $\sum_{i\in\mc I}Fi$ is an equivalence relation,
  \[ (i,x) \sim (j,y) \coloneq \ex{k}{\mc I}\ex{\alpha}{\mc I(i,k)}\ex{\beta}{\mc I(j,k)} F\alpha(x) = F\beta(y). \]
\end{lemma}
\begin{proof}
  Reflexivity and symmetry hold by construction. Transitivity holds by the fact that any $\kappa$-filtered category admits finite cocones.
\end{proof}

This leads to the following well-known result:

\begin{theorem}\label{thm:filtercommute}
  $\kappa$-filtered colimits commutes with $\kappa$-small limits in $\Set$: for any diagram $F \colon \mc I \times \mc D \to \Set$ where $\mc I$ is $\kappa$-filtered and $\mc D$ is $\kappa$-small, the canonical map below is an isomorphism,
  \[ \ct_{i\in\mc I}\lt_{d\in\mc D} F(i,d) \to \lt_{d\in\mc D}\ct_{i\in\mc I}F(i,d). \]
\end{theorem}
\begin{proof}
  The map sends the equivalence class $[(i,x)]$ to the family $\ld{d}{\mc D}[(i,x_d)]$. To show it is injective, we may assume we have $[(i,x)]$ and $[(j,y)]$ agreeing in the target,
  \[ \fa{d}{\mc D}[(i,x_d)] = [(j,y_d)]. \]
  By Lemma~\ref{lem:filteredquotient}, this explicitly means that
  \[ \fa{d}{\mc D}\ex{k}{\mc I}\ex{\alpha}{\mc I(i,k)}\ex{\beta}{\mc I(j,k)} F(d,\alpha)(x_d) = F(d,\beta)(y_d). \]
  Since $\mc D$ is $\kappa$-small, its set of objects is projective. Hence we can assume that there exists a function $k \colon \Ob(\mc D) \to \Ob(\mc I)$ and similarly $\alpha,\beta$ picking out the morphisms satisfying the above condition. The resulting diagram on $\mc I$ will be $\kappa$-small, thus by assumption we can pick a cocone of it, which we denote as $(u,\gamma)$. In particular, it satisfies that
  \[ \fa{d,d'}{\mc D} \gamma_d \circ \alpha_d = \gamma_{d'} \circ \alpha_{d'} \conjt \gamma_d \circ \beta_d = \gamma_{d'} \circ \beta_{d'}. \]
  Now we have $u$ with $\gamma\circ\alpha \in \mc I(i,u)$ and $\gamma \circ \beta \in \mc I(j,u)$. To show $F(\gamma\circ\alpha)(x) = F(\gamma\circ\beta)(y)$, it suffices to show
  \[  \fa{d}{\mc D}F(d,\gamma_d\circ\alpha_d)(x_d) = F(d,\gamma_d\circ\beta_d)(y_d), \]
  which holds by construction that $(u,\gamma)$ is a cocone. Hence, $[(i,x)] = [(j,y)]$.

  For surjectivity, suppose we have a term in $\prod_{d\in\mc D}\ct_{iin\mc I}F(i,d)$. Again by $\Ob(\mc D)$ being projective, we may assume we have chosen a representative $[(i_d,x_d)]$ for all $d \in \mc D$, satisfying for any $\alpha \in \mc D(d,e)$,
  \[ [(i_d,F(i_d,\alpha)(x_d))] = [(i_e,x_e)]. \]
  In particular, this means there is $j_{d,e,\alpha} \in \mc I$, $\beta_{d,e,\alpha} \in \mc I(i_d,j_{d,e,\alpha})$, and $\gamma_{d,e,\alpha} \in \mc I(i_e,j_{d,e,\alpha})$, that
  \[ F(\beta_{d,e,\alpha},\alpha)(x_d) = F(\gamma_{d,e,\alpha},e)(x_e). \]
  Now the diagram indexed by $\sum_{d,e\in\mc D}\mc D(d,e)$ is again $\kappa$-small, hence we can find a cocone $k$ with $\delta,\xi$, such that for any $d,e \in \mc D_0$ and $\alpha \in \mc D(d,e)$,
  \[ \delta_{d,e,\alpha} \circ \beta_{d,e,\alpha} = \xi_d, \quad \delta_{d,e,\beta} \circ \gamma_{d,e,\alpha} = \xi_e. \]
  This way, we can pick the representative $(k,x)$ with
  \[ x \coloneq \ld{d}{\mc D} F(\xi_d,d)(x_d). \]
  This is indeed well-defined, because for any $\alpha \in \mc D(d,e)$,
  \[ F(\xi_d,\alpha)(x_d) = F(\delta_{d,e,\alpha}\circ\beta_{d,e,\alpha},\alpha)(x_d) = F(\delta_{d,e,\alpha}\circ\gamma_{d,e,\alpha},e)(x_e) = F(\xi_e,e)(x_e). \]
  Thus, we only need to observe that $[(k,d)]$ is mapped to $\ld d{\mc D}[(i_d,x_d)]$, which holds by construction.
\end{proof}



\section{Locally presentable categories}

Again we fix a separated notion of size $\kappa$.

\begin{definition}[Presentable objects]
  Let $\mc A$ be a category having $\kappa$-filtered colimits. An object $a \in \mc A$ is \emph{$\kappa$-presentable}, if the representable functor $\mc A(a,-) \colon \mc A \to \Set$ preserves $\kappa$-filtered colimits. The full subcategory of $\kappa$-presentable objects will be denoted as $\mc A_\kappa$.
\end{definition}

\begin{remark}[Finitely and countably presentable sets]
  We will see later in~\zcref{exa:presentableobjectsinset} that $\Set_{\omega} = \Fin$. On the other hand, constructively $\Set_{\omega_1}$ may not coincide with $\Cnt$, since the former will be closed under countable colimits (as we will soon see in~\zcref{lem:presentableclosedundercolimits}), in particular coequalisers, while as we have seen in~\zcref{rem:coequalisercntandmarkov} the latter may not be.
\end{remark}

\begin{lemma}\label{lem:presentableclosedundercolimits}
  $\mc A_{\kappa}$ is closed under $\kappa$-small colimits that exists in $\mc A$.
\end{lemma}
\begin{proof}
  This is a formal consequence of Theorem~\ref{thm:filtercommute}.
\end{proof}

To define the notion of locally $\kappa$-presentable categories, we first consider the following notion of \emph{generators} (cf.~\citet[Def. 4.5.4]{borceux1994handbookI}):

\begin{definition}[Dense subcategory]
  For any category $\mc A$, a full subcategory $\mc B \subseteq \mc A$ is \emph{dense} if for any $a \in \mc A$, the canonical map below is an isomorphism,
  \[ \ct_{b\in\mc B/a}b \to a. \]
\end{definition}

There is an equivalent criterion for dense subcategories:

\begin{lemma}\label{lem:densefullyfaithful}
  $\mc B \subseteq \mc A$ is dense iff the restricted Yoneda functor below is fully faithful,
  \[ R \colon \mc A \to \psh(\mc B). \]
\end{lemma}
\begin{proof}
  See~\citet[Prop. 4.5.14]{borceux1994handbookI}; the proof there is constructive.
\end{proof}

\begin{definition}[Locally presentable categories]\label{def:locallypresentablecategory}
  A category $\mc A$ is \emph{locally $\kappa$-presentable}, if it is cocomplete and $\mc A_\kappa$ is dense and essentially small.
\end{definition}

\begin{remark}[Locally presentable categories are locally small]
  If $\mc A$ is locally $\kappa$-presentable then it must be locally small, because it is a full subcategory of $\psh(\mc A_\kappa)$ for the essentially small category $\mc A_\kappa$.
\end{remark}

\begin{example}[Presheaf categories are locally $\omega$-presentable]\label{exm:pshupresentable}
  For any small category $\mc B$, the presheaf category $\psh(\mc B)$ is locally $\omega$-presentable.
\end{example}

In general, we have the following fact:

\begin{proposition}\label{prop:locallypresaccembedpsh}
  Let $\mc A$ be locally $\kappa$-presentable. The restricted Yoneda embedding
  \[ R \colon \mc A \hook \psh(\mc A_\kappa), \]
  preserves $\kappa$-filtered colimits, and it has a left adjoint $L$.
\end{proposition}
\begin{proof}
  For any $\kappa$-filtered colimit $a \cong \ct_{i\in\mc I}a_i$ in $\mc A$ and $b \in \mc A_\kappa$,
  \[ Ra(b) \cong \mc A(b,a) \cong \ct_{i\in\mc I}\mc A(b,a_i) \cong \ct_{i\in\mc I}Ra_i(b). \]
  This uses the fact that $b$ is $\kappa$-presentable, hence it follows that $R$ preserves $\kappa$-filtered colimits. It is also fully faithful because $\mc A_\kappa$ is dense. The left adjoint is simply given by point-wise left Kan extension, which takes a presheaf $F$ to the following colimit,
  \[ LF \coloneq \ct_{b\in\mc A_\kappa/F}b. \]
  It is indeed a left adjoint because we have
  \begin{align*}
    &\mc A(LF,a) \\
    &\quad\cong \lt_{b\in\mc A_\kappa/F}\mc A(b,a) \\
    &\quad\cong \lt_{b\in\mc A_\kappa/F}\ct_{c\in\mc A_\kappa/a}\mc A(b,c) \\
    &\quad\cong \lt_{b\in\mc A_\kappa/F}\ct_{c\in\mc A_\kappa/a}\psh(\mc A_\kappa)(\yon_b,\yon_c) \\
    &\quad\cong \psh(\mc A_\kappa)(\ct_{b\in\mc A_\kappa/F}\yon_b,\ct_{c\in\mc A_\kappa/a}\yon_c) \\
    &\quad\cong \psh(\mc A_\kappa)(F,Ra)
  \end{align*}
  The first isomorphism holds by definition of $F$; the second holds since $b$ is $\kappa$-presentable, and that $\mc A_\kappa/a$ is $\kappa$-filtered by~\zcref{lem:presentableclosedundercolimits} and~\zcref{rem:kappacompeltefiltered}; the third holds by Yoneda; the forth holds since representables preserve all colimits; the last holds by definition and the fact that $R$ preserves $\kappa$-filtered colimits.
\end{proof}

\begin{corollary}
  If $\mc A$ is locally $\kappa$-presentable, then it is complete.
\end{corollary}
\begin{proof}
  By Proposition~\ref{prop:locallypresaccembedpsh}, $\mc A$ is a reflective subcategory of a presheaf category.
\end{proof}

Let $\kappa\hp\Lex$ denote the 2-category of categories with $\kappa$-small limits and functors preserving $\kappa$-small limits. We can characterise the essential image of the embedding $\mc A \hook \psh(\mc A_\kappa)$ for a locally $\kappa$-presentable category $\mc A$ as follows:

\begin{proposition}\label{prop:presentableindlexfunctor}
  For any locally $\kappa$-presentable category $\mc A$, 
  \[ \mc A \simeq \kappa\hp\Lex(\mc A_{\kappa}\op,\Set). \]
\end{proposition}
\begin{proof}
  For any $a\in\mc A$, by construction $Ra \cong \mc A(-,a)$ takes any colimit in $\mc A_\kappa$ to a limit in $\Set$, thus $R$ factors through $\kappa\hp\Lex(\mc A_\kappa\op,\Set)$. On the other hand, if $F \colon \mc A_{\kappa}\op \to \Set$ preserves $\kappa$-small limits, then its category of elements $\mc A_\kappa/F$ will be $\kappa$-filtered; a simple proof can be given through the characterisation of $\kappa$-filtered diagrams in~\zcref{lem:concretefilter}. This way,
  \[ F \cong \ct_{a\in\mc A_\kappa/F}\yon_a \cong R\ct_{a\in\mc A_\kappa/F}a, \]
  where the second isomorphism holds by the fact that $R$ preserves $\kappa$-filtered colimits by~\zcref{prop:locallypresaccembedpsh}. This shows the desired equivalence holds.
\end{proof}

% \begin{proposition}\label{prop:upresencloureflction}
%   If $\mc A$ is locally $\kappa$-presentable, and $\mc B$ is a full reflective subcategory of $\mc A$ closed under $\kappa$-filtered colimits, then $\mc B$ is also locally $\kappa$-presentable.
% \end{proposition}
% \begin{proof}
%   Let us use $R \colon \mc B \to \mc A$ to denote the inclusion and $L$ the reflection. Note that $\mc B$ is cocomplete, and for any $F\colon\mc I \to \mc B$, we have
%   \[ \ct_{u\in\mc I}Fu = L\ct_{u\in\mc I}RFu. \]
%   This means that if $i \in \mc A_{\kappa}$ then $Li \in \mc B_{\kappa}$, because for any $\kappa$-filtered diagram $F \colon \mc I \to \mc B$,
%   \begin{align*}
%     &\mc B(Li,\ct_{u\in\mc I}Fu) \\
%     &\quad\cong \mc A(i,R\ct_{u\in\mc I}Fu) \\
%     &\quad\cong \mc A(i,\ct_{u\in\mc I}RFu) \\
%     &\quad\cong \ct_{u\in\mc I}\mc A(i,RFu) \\
%     &\quad\cong \ct_{u\in\mc I}\mc B(Li,Fu).
%   \end{align*}
%   The second step uses the fact that $R$ preserves $\kappa$-filtered colimits; the other steps are routine. Hence, $\mc B_{\kappa}$ is the image of $\mc A_{\kappa}$ under $L$. 

%   Since we are proving a proposition, we can assume there is a small dense full subcategory $\mc U$ of $\kappa$-presentable objects in $\mc A$. Let $\mc V$ be the full subcategory of $\mc B$ under the image of $L$ of $\mc U$. Now for any $b : \mc B_0$, we have
%   \[ \ct_{i:\mc V/b}i = L\ct_{i:\mc V/b}Ri. \]
%   Note that there is a canonical map
%   \[ Rb = \ct_{j:\mc U/Rb}j \to \ct_{i:\mc V/b}Ri \]
%   induced by the functor $L : \mc U/Rb \to \mc V/b$ and the unit of the adjunction $L \dashv R$. This map is indeed localised by $L$, because the unit $j \to RLj$ are localised by $L$. This way, we have
%   \[ b = LRb = L\ct_{j:\mc U/Rb}j = L\ct_{i:\mc V/b}Ri = \ct_{i:\mc V/b}LRi = \ct_{i:\mc V/b}i. \]
%   Hence, $\mc B$ is locally $\kappa$-presentable.
% \end{proof}

% \begin{theorem}
%   The following are equivalent for a category $\mc A$:
%   \begin{enumerate}
%     \item $\mc A$ is locally $\kappa$-presentable;
%     \item $\mc A$ is a full reflective subcategory of $\psh(\mc B)$ for a small category $\mc B$, which are closed under $\kappa$-filtered colimits.
%   \end{enumerate}
% \end{theorem}
% \begin{proof}
%   (1) $\nt$ (2) is Proposition~\ref{prop:locallypresaccembedpsh}. For (2) $\nt$ (1), it follows by Example~\ref{exm:pshupresentable} and Proposition~\ref{prop:upresencloureflction}.
% \end{proof}